\documentclass[11pt]{article}
\usepackage[a4paper,margin=22mm]{geometry}
\usepackage{fontspec}
\setmainfont{DejaVu Serif}
\setsansfont{DejaVu Sans}
\usepackage{microtype}
\usepackage{xcolor}
\usepackage{hyperref}
\usepackage{enumitem}
\usepackage{parskip}
\definecolor{Accent}{HTML}{971D32}
\definecolor{Muted}{HTML}{666666}
\hypersetup{colorlinks=true,urlcolor=Accent,linkcolor=Accent}
\setlist{leftmargin=1.6em,itemsep=0.3em,topsep=0.35em}
\setcounter{tocdepth}{2}
\renewcommand{\contentsname}{Contents}
\newcommand{\sectionrule}{\par\vspace{-0.5em}\noindent\textcolor{Accent}{\rule{\linewidth}{0.6pt}}\par}
\begin{document}

\begin{center}
{\sffamily\bfseries\color{Accent} TOEFL WORD OF THE DAY\par}
\vspace{8mm}
{\fontsize{42}{48}\selectfont\bfseries attribute\par}
\vspace{2mm}
{\Large\sffamily\color{Accent} /ˈætrəˌbjuːt/\par}
\vspace{2mm}
{\small\sffamily Local audio phrase: an attribute\par}
{\small\sffamily Audio file: audios/an-attribute.mp3\par}
\vspace{3mm}
{\bfseries\color{Accent} Pronunciation changes across meanings.\par}
countable noun: /ˈætrəˌbjuːt/ --- “an attribute”\par
transitive verb: /əˈtrɪbjuːt/ --- “attribute the change”\par
\vspace{2mm}
{\large\sffamily\color{Muted} countable noun and transitive verb\par}
\vspace{5mm}
{\sffamily an important quality or feature \textbullet\ regard something as caused by \textbullet\ regard something as created by\par}
\vspace{6mm}
{\large An attribute is a quality or feature, while the verb attribute connects a result, work, or statement with its believed cause or creator.\par}
\vspace{6mm}
\begin{minipage}{0.84\textwidth}
\itshape “Scientists attribute the decline in the bird population to habitat loss.”
\end{minipage}\par
\vspace{10mm}
\textbf{Date:} August 10th 2026\quad\textbullet\quad \textbf{Level:} B1--mid-B2 TOEFL
\vfill
Created and compiled by \textbf{Amir Kooshky}\par
\vspace{2mm}
\href{https://t.me/KooshkyTOEFL}{Telegram Channel}\quad\textbullet\quad
\href{https://www.instagram.com/kooshkytoefl}{Instagram}\quad\textbullet\quad
\href{https://t.me/KooshkyTOEFL\_pv}{Personal Telegram}
\end{center}
\newpage
\tableofcontents
\newpage

\section{Meanings and usage}
\sectionrule
\subsection{Meaning 1: Quality or feature}
\textbf{Part of speech:} countable noun

\textbf{IPA:} /ˈætrəˌbjuːt/

\textbf{Pronunciation phrase:} an attribute

\textbf{Local audio file:} audios/an-attribute.mp3

\textbf{Register:} neutral to formal; common in academic, professional, scientific, and technical writing

\textbf{Definition:} a quality or feature that is considered a typical or important part of a person, thing, or system

\textbf{In simpler words:} an important quality or feature

\textbf{Typical pattern:} an attribute of + person, object, or system

\subsubsection{Natural examples}
\begin{enumerate}
  \item Patience is an important attribute of an effective teacher.
  \item Flexibility is one of the material's most useful attributes.
  \item The researchers examined several attributes of successful organizations.
  \item Physical strength was once regarded as an essential attribute of military leadership.
  \item The ability to adapt is a valuable attribute in a rapidly changing workplace.
  \item Each plant species has attributes that help it survive in a particular environment.
  \item Accuracy and clarity are desirable attributes in academic writing.
  \item The software allows users to search for products according to specific attributes.
\end{enumerate}

\subsubsection{Nuance and implication}
\textbf{Central nuance:} An attribute is a quality that belongs to or helps describe a person, object, system, or category.

\textbf{Usual implication:} The quality is often considered characteristic, important, or useful for identifying or evaluating something.

\textbf{Compared with possession:} A possession is something that a person or organization owns. An attribute is a quality or feature that something has.

\subsubsection{Grammar notes}
\textbf{How to use it:} Use an attribute for one quality and attributes for several qualities.

\textbf{What can follow it:} It is often followed by of and the person, object, or system that has the quality.

\textbf{Common preposition:} Use an attribute of something, as in an important attribute of good leadership.

\textbf{Noun use:} This is a countable noun, so the singular normally needs a, an, or the.

\subsubsection{Useful patterns}
\textbf{Pattern:} an attribute of + person, thing, or system

\textbf{Use:} Use this to identify something that possesses a particular quality.

\textbf{Example:} Reliability is a central attribute of a well-designed transportation system.

\textbf{Pattern:} a desirable or valuable attribute

\textbf{Use:} Use this for a quality that is considered useful or positive.

\textbf{Example:} Curiosity is a valuable attribute in scientific research.

\textbf{Pattern:} a defining attribute

\textbf{Use:} Use this for a quality that strongly distinguishes something from other things.

\textbf{Example:} The ability to regulate body temperature is a defining attribute of mammals.

\textbf{Pattern:} possess or have an attribute

\textbf{Use:} Use this when stating that a person or thing has a particular quality.

\textbf{Example:} The material possesses several attributes that make it suitable for construction.

\subsubsection{Common wrong usage}
\paragraph{Correction 1}
\textbf{Wrong:} Patience is an attribute for a good teacher.

\textbf{Correct:} Patience is an attribute of a good teacher.

\textbf{Why:} Use an attribute of someone or something.

\paragraph{Correction 2}
\textbf{Wrong:} Honesty is important attribute in a leader.

\textbf{Correct:} Honesty is an important attribute in a leader.

\textbf{Why:} The singular countable noun attribute needs an article.

\paragraph{Correction 3}
\textbf{Wrong:} The machine has many useful attributions.

\textbf{Correct:} The machine has many useful attributes.

\textbf{Why:} Attributes are qualities or features. Attributions are explanations of causes or statements about who created something.

\subsection{Meaning 2: Regard as caused by}
\textbf{Part of speech:} transitive verb

\textbf{IPA:} /əˈtrɪbjuːt/

\textbf{Pronunciation phrase:} attribute the change

\textbf{Local audio file:} audios/attribute-the-change.mp3

\textbf{Register:} formal; very common in academic, scientific, historical, medical, and professional writing

\textbf{Definition:} to believe that an event, condition, or result was caused by a particular person, thing, or factor

\textbf{In simpler words:} believe that something caused a result

\textbf{Typical pattern:} attribute + result or condition + to + cause

\subsubsection{Natural examples}
\begin{enumerate}
  \item Scientists attribute the decline in the bird population to habitat loss.
  \item She attributed her success to careful preparation and support from her colleagues.
  \item The improvement cannot be attributed entirely to the new policy.
  \item Researchers attributed the difference between the groups to variations in diet.
  \item Some historians attribute the collapse of the society to a prolonged drought.
  \item The company attributed the delay to unexpected technical problems.
  \item Doctors initially attributed his symptoms to stress.
  \item Analysts attribute much of the recent economic growth to increased consumer spending.
\end{enumerate}

\subsubsection{Nuance and implication}
\textbf{Central nuance:} The verb identifies what someone believes to be the cause or explanation of a result.

\textbf{Usual implication:} The connection may be supported by evidence, but attribute does not necessarily mean that the cause has been proved with complete certainty.

\textbf{Compared with prove:} To attribute a result to something is to identify it as the believed cause. To prove that it caused the result requires stronger and more conclusive evidence.

\subsubsection{Grammar notes}
\textbf{How to use it:} Place the result, event, or condition directly after attribute.

\textbf{What can follow it:} After the result, use to and then the believed cause.

\textbf{Common preposition:} Use attribute something to something, as in attribute the decline to pollution.

\textbf{Position in a sentence:} The passive form is very common: The decline was attributed to pollution.

\subsubsection{Useful patterns}
\textbf{Pattern:} attribute + result + to + cause

\textbf{Use:} Use this when identifying the believed cause of an event or condition.

\textbf{Example:} Researchers attributed the increase in crop production to improved irrigation.

\textbf{Pattern:} be attributed to + cause

\textbf{Use:} Use the passive form when the result is more important than the person explaining it.

\textbf{Example:} The population decline was attributed to disease and food shortages.

\textbf{Pattern:} attribute success or failure to + factor

\textbf{Use:} Use this when explaining why a person, project, or organization succeeded or failed.

\textbf{Example:} The team attributed its success to careful planning.

\textbf{Pattern:} largely, partly, or mainly attribute something to + cause

\textbf{Use:} Use these words to show how much of the result is explained by a particular factor.

\textbf{Example:} Economists largely attribute the recovery to increased investment.

\subsubsection{Common wrong usage}
\paragraph{Correction 1}
\textbf{Wrong:} Researchers attribute the decline from climate change.

\textbf{Correct:} Researchers attribute the decline to climate change.

\textbf{Why:} Use attribute something to a cause.

\paragraph{Correction 2}
\textbf{Wrong:} The company attributed to poor planning the failure.

\textbf{Correct:} The company attributed the failure to poor planning.

\textbf{Why:} The usual order is attribute the result to the cause.

\paragraph{Correction 3}
\textbf{Wrong:} The improvement was attributed by the new treatment.

\textbf{Correct:} The improvement was attributed to the new treatment.

\textbf{Why:} Use to before the believed cause. By would identify the person making the attribution.

\paragraph{Correction 4}
\textbf{Wrong:} Scientists attributed that pollution caused the decline.

\textbf{Correct:} Scientists attributed the decline to pollution.

\textbf{Why:} Use attribute a result to a cause rather than attribute that followed by a clause.

\subsection{Meaning 3: Regard as created or said by}
\textbf{Part of speech:} transitive verb

\textbf{IPA:} /əˈtrɪbjuːt/

\textbf{Pronunciation phrase:} attribute the painting

\textbf{Local audio file:} audios/attribute-the-painting.mp3

\textbf{Register:} formal; common in history, art, literature, science, and academic writing

\textbf{Definition:} to believe or state that a work, statement, idea, or achievement was created, said, or accomplished by a particular person or group

\textbf{In simpler words:} believe that someone created or said something

\textbf{Typical pattern:} attribute + work, quotation, discovery, or idea + to + person or group

\subsubsection{Natural examples}
\begin{enumerate}
  \item The painting was originally attributed to a famous Italian artist.
  \item This quotation is often attributed to Einstein, although he may never have said it.
  \item Historians attribute the invention to several different researchers.
  \item The poem was wrongly attributed to a writer who lived a century later.
  \item Scholars have attributed the ancient manuscript to an unknown religious community.
  \item The discovery is generally attributed to a team of French scientists.
  \item Several architectural innovations are attributed to the Roman engineer.
  \item The theory was once attributed to Darwin but was actually proposed by an earlier scholar.
\end{enumerate}

\subsubsection{Nuance and implication}
\textbf{Central nuance:} The verb connects a work, statement, idea, or achievement with the person believed to be responsible for it.

\textbf{Usual implication:} The identification may be accepted, uncertain, disputed, or later shown to be incorrect.

\textbf{Compared with credit:} To credit someone with something usually emphasizes recognizing their contribution or achievement. To attribute something to someone focuses on identifying them as its source, creator, or author.

\subsubsection{Grammar notes}
\textbf{How to use it:} Place the work, statement, idea, or achievement directly after attribute.

\textbf{What can follow it:} Use to and then the person, group, culture, or period believed to be responsible.

\textbf{Common preposition:} Use attribute something to someone, as in attribute the poem to an unknown writer.

\textbf{Position in a sentence:} The passive form is especially common: The quotation is attributed to Mark Twain.

\subsubsection{Useful patterns}
\textbf{Pattern:} attribute + work or statement + to + person

\textbf{Use:} Use this when identifying the believed author, speaker, or creator.

\textbf{Example:} Experts attributed the sculpture to a local artist.

\textbf{Pattern:} be attributed to + person or group

\textbf{Use:} Use the passive form when the work or achievement is the main focus.

\textbf{Example:} The discovery is usually attributed to two independent research teams.

\textbf{Pattern:} wrongly or mistakenly attribute something to + person

\textbf{Use:} Use this when the accepted source or creator is incorrect.

\textbf{Example:} The quotation is frequently misattributed to Abraham Lincoln.

\textbf{Pattern:} generally or traditionally attribute something to + person

\textbf{Use:} Use this when an attribution is widely accepted but may not be completely certain.

\textbf{Example:} The epic poem is traditionally attributed to Homer.

\subsubsection{Common wrong usage}
\paragraph{Correction 1}
\textbf{Wrong:} The quotation is attributed from a famous scientist.

\textbf{Correct:} The quotation is attributed to a famous scientist.

\textbf{Why:} Use attribute something to a person.

\paragraph{Correction 2}
\textbf{Wrong:} Historians attribute the scientist for the discovery.

\textbf{Correct:} Historians attribute the discovery to the scientist.

\textbf{Why:} Place the discovery after attribute and the person after to.

\paragraph{Correction 3}
\textbf{Wrong:} The painting was attributed as Leonardo da Vinci.

\textbf{Correct:} The painting was attributed to Leonardo da Vinci.

\textbf{Why:} Use to, not as, before the believed creator.

\section{Word family}
\sectionrule
\subsection{attribution}
\textbf{Part of speech:} countable or uncountable noun

\textbf{IPA:} /ˌætrəˈbjuːʃən/

\textbf{Definition:} the act of explaining something as the result of a particular cause, or of identifying someone as the creator of a work or idea

\textbf{Relationship to the headword:} This noun names the act or result of attributing a cause, work, statement, or achievement.

\textbf{Grammar pattern:} the attribution of + result + to + cause or the attribution of + work + to + creator

\textbf{Usage note:} It can be countable when referring to a particular explanation or identification and uncountable when referring to the general practice.

\subsubsection{Examples}
\begin{enumerate}
  \item The attribution of the decline to climate change remains controversial.
  \item Experts questioned the painting's attribution to the famous artist.
  \item Accurate attribution is essential when using another person's ideas.
\end{enumerate}

\subsection{attributable}
\textbf{Part of speech:} adjective

\textbf{IPA:} /əˈtrɪbjətəbəl/

\textbf{Definition:} considered to have been caused by a particular person, thing, or factor

\textbf{Relationship to the headword:} This adjective describes a result that can be attributed to a particular cause.

\textbf{Grammar pattern:} be attributable to + cause

\textbf{Usage note:} It is formal and especially common in academic, legal, scientific, and financial writing.

\subsubsection{Examples}
\begin{enumerate}
  \item Most of the improvement was attributable to better medical care.
  \item The damage attributable to the storm was less severe than expected.
\end{enumerate}

\section{Common collocations}
\sectionrule
\subsection{important attribute}
\textbf{Applies to:} Quality or feature

\textbf{Meaning:} a quality that has considerable value or significance

\textbf{Pattern:} an important attribute of + person, object, or system

\textbf{Example:} The ability to cooperate is an important attribute of effective leadership.

\subsection{key attribute}
\textbf{Applies to:} Quality or feature

\textbf{Meaning:} one of the most important qualities of something

\textbf{Example:} Durability is a key attribute of the new material.

\subsection{defining attribute}
\textbf{Applies to:} Quality or feature

\textbf{Meaning:} a quality that strongly identifies or distinguishes something

\textbf{Example:} Complex language is often considered a defining attribute of human societies.

\subsection{personal attribute}
\textbf{Applies to:} Quality or feature

\textbf{Meaning:} a quality belonging to an individual person

\textbf{Example:} Employers often value personal attributes such as reliability and adaptability.

\subsection{attribute a change to}
\textbf{Applies to:} Regard as caused by

\textbf{Meaning:} identify a particular factor as the cause of a change

\textbf{Pattern:} attribute a change to + cause

\textbf{Example:} Researchers attributed the change in behavior to increased social contact.

\subsection{attribute success to}
\textbf{Applies to:} Regard as caused by

\textbf{Meaning:} identify the factors believed to have produced success

\textbf{Pattern:} attribute success to + factor

\textbf{Example:} She attributes her success to persistence and careful planning.

\subsection{largely attributable to}
\textbf{Applies to:} Regard as caused by

\textbf{Meaning:} mainly considered to have been caused by something

\textbf{Pattern:} be largely attributable to + cause

\textbf{Example:} The reduction in mortality was largely attributable to improved sanitation.

\textbf{Usage note:} Attributable is the adjective form.

\subsection{directly attributable to}
\textbf{Applies to:} Regard as caused by

\textbf{Meaning:} considered to be a direct result of something

\textbf{Pattern:} be directly attributable to + cause

\textbf{Example:} Only a small part of the damage was directly attributable to the new policy.

\textbf{Usage note:} Attributable is the adjective form.

\subsection{attribute a quotation to}
\textbf{Applies to:} Regard as created or said by

\textbf{Meaning:} identify someone as the believed speaker or writer of a quotation

\textbf{Pattern:} attribute a quotation to + person

\textbf{Example:} Many websites attribute the quotation to a writer who never used those exact words.

\subsection{attribute a work to}
\textbf{Applies to:} Regard as created or said by

\textbf{Meaning:} identify someone as the believed creator of a work

\textbf{Pattern:} attribute a work to + creator

\textbf{Example:} Art historians attributed the work to one of the master's students.

\subsection{widely attributed to}
\textbf{Applies to:} Regard as created or said by

\textbf{Meaning:} generally believed to have been created or accomplished by someone

\textbf{Example:} The invention is widely attributed to a nineteenth-century engineer.

\subsection{mistakenly attributed to}
\textbf{Applies to:} Regard as created or said by

\textbf{Meaning:} incorrectly identified as the work or statement of someone

\textbf{Example:} The statement is often mistakenly attributed to the former president.

\section{Synonyms and antonyms}
\sectionrule
\subsection{Close synonyms}
\subsubsection{characteristic}
\textbf{Part of speech:} countable noun

\textbf{Applies to:} Quality or feature

\textbf{Shared meaning:} Both words can name a quality or feature of a person or thing.

\textbf{Difference:} Characteristic is a common general word for a typical quality. Attribute is slightly more formal and may refer to any important or identifiable quality, whether or not it is typical.

\textbf{Example:} Thick fur is a characteristic of many animals that live in cold climates.

\subsubsection{trait}
\textbf{Part of speech:} countable noun

\textbf{Applies to:} Quality or feature

\textbf{Shared meaning:} Both words can refer to a quality that helps describe someone or something.

\textbf{Difference:} Trait is especially common for qualities of a person's character or inherited features of an organism. Attribute can describe people, objects, systems, and abstract concepts.

\textbf{Example:} Curiosity is a valuable personality trait.

\subsubsection{ascribe}
\textbf{Part of speech:} transitive verb

\textbf{Applies to:} all senses

\textbf{Shared meaning:} Both verbs can connect a result with a cause or a work with its believed creator.

\textbf{Difference:} Ascribe is a close formal synonym of the verb attribute. Attribute is more common, especially in academic and scientific writing.

\textbf{Example:} The researchers ascribed the improvement to changes in diet.

\subsubsection{credit}
\textbf{Part of speech:} transitive verb

\textbf{Applies to:} Regard as created or said by

\textbf{Shared meaning:} Both verbs can identify someone as responsible for an achievement or idea.

\textbf{Difference:} Credit often emphasizes giving recognition or praise for an achievement. Attribute focuses more neutrally on identifying the believed source or creator.

\textbf{Example:} The team credited its coach with developing the successful strategy.

\section{Words learners may confuse}
\sectionrule
\subsection{contribute}
\textbf{Part of speech:} transitive verb versus intransitive or transitive verb

\textbf{Why learners confuse them:} Both words frequently appear in explanations of causes, but the direction of the relationship is different.

\textbf{Key difference:} Attribute means connect a result with its believed cause. Contribute means help cause a result or give something toward a shared effort.

\textbf{attribute:} Researchers attribute the decline to habitat destruction.

\textbf{contribute:} Habitat destruction contributed to the decline.

\subsection{attribution}
\textbf{Part of speech:} noun or verb versus noun

\textbf{Why learners confuse them:} The words share the same root but perform different grammatical jobs.

\textbf{Key difference:} Attribute can be a noun meaning a quality or a verb meaning connect something with a cause or creator. Attribution is the noun for the act or result of making that connection.

\textbf{attribute:} Flexibility is a useful attribute of the material.

\textbf{attribution:} The attribution of the painting to that artist remains uncertain.

\section{Etymology}
\sectionrule
\textbf{Origin:} Attribute ultimately comes from Latin attribuere, meaning to assign or give to.

\section{Practice exercises}
\sectionrule
\subsection{Word family choice}
Choose the best member of the word family for each blank.

\begin{enumerate}
  \item The decline was largely \_\_\_\_\_ to the destruction of natural habitats.
  \item Experts questioned the \_\_\_\_\_ of the painting to the famous artist.
  \item Adaptability is a valuable \_\_\_\_\_ in a rapidly changing workplace.
\end{enumerate}

\subsection{Correct or incorrect?}
Decide whether each sentence is correct. Correct the inaccurate ones.

\begin{enumerate}
  \item Researchers attributed the increase in crop production to improved irrigation.
  \item Patience is an important attribute for an effective teacher.
  \item The quotation is often attributed from a famous scientist.
  \item The improvement was partly attributable to better training.
  \item The researchers attributed that the treatment caused the improvement.
\end{enumerate}

\subsection{Collocation practice}
Complete each sentence with a natural collocation.

\begin{enumerate}
  \item The researchers attribute the decline \_\_\_\_\_ changes in the local climate.
  \item Honesty is a key \_\_\_\_\_ of effective leadership.
  \item The ancient text is traditionally attributed \_\_\_\_\_ an unknown poet.
  \item The rapid improvement was largely \_\_\_\_\_ to increased investment.
\end{enumerate}

\subsection{Complete answer key}
\subsubsection{Word family choice}
\begin{enumerate}
  \item \textbf{attributable.} The decline was largely attributable to the destruction of natural habitats. \emph{Use the adjective attributable after was and before to.}
  \item \textbf{attribution.} Experts questioned the attribution of the painting to the famous artist. \emph{Use the noun attribution for the identification of the painting's believed creator.}
  \item \textbf{attribute.} Adaptability is a valuable attribute in a rapidly changing workplace. \emph{Use the noun attribute for a quality or feature.}
\end{enumerate}

\subsubsection{Correct or incorrect?}
\begin{enumerate}
  \item \textbf{Correct} \emph{Attribute a result to a cause is the standard pattern.}
  \item \textbf{Incorrect}. Patience is an important attribute of an effective teacher. \emph{Use an attribute of a person or thing.}
  \item \textbf{Incorrect}. The quotation is often attributed to a famous scientist. \emph{Use attribute something to someone.}
  \item \textbf{Correct} \emph{Be attributable to correctly means be considered partly caused by something.}
  \item \textbf{Incorrect}. The researchers attributed the improvement to the treatment. \emph{Use attribute the result to the cause.}
\end{enumerate}

\subsubsection{Collocation practice}
\begin{enumerate}
  \item \textbf{to.} The researchers attribute the decline to changes in the local climate. \emph{Use attribute something to a cause.}
  \item \textbf{attribute.} Honesty is a key attribute of effective leadership. \emph{A key attribute is an important quality or feature.}
  \item \textbf{to.} The ancient text is traditionally attributed to an unknown poet. \emph{Use be attributed to when identifying the believed creator.}
  \item \textbf{attributable.} The rapid improvement was largely attributable to increased investment. \emph{Largely attributable to means mainly considered to have been caused by something.}
\end{enumerate}

\vfill
\begin{center}
\small\color{Muted} Created and compiled by \textbf{Amir Kooshky}\\
\href{https://t.me/KooshkyTOEFL}{Telegram Channel}\quad\textbullet\quad
\href{https://www.instagram.com/kooshkytoefl}{Instagram}\quad\textbullet\quad
\href{https://t.me/KooshkyTOEFL\_pv}{Personal Telegram}
\end{center}
\end{document}
